\section{Asymptotic expansion, explicit remainders, order, and type}\label{sec:asymptotic}

The symbol \(\sim\) in an asymptotic series must be replaced in iPRA by a
family of finite remainder inequalities.  For Barnes $G$, the most convenient
starting point is the Ferreira--L\'opez form with Nemes' explicit remainder.

\subsection{The Gamma-coupled expansion}

For an integer \(N\ge1\), define
\begin{align}
 \mathcal B_N(z)={}&\frac14z^2+z\Log\Gamma(z+1)
 -\left(\frac12z(z+1)+\frac1{12}\right)\Log z-\log A\notag\\
 &+\sum_{k=1}^{N-1}
 \frac{B_{2k+2}}
 {2k(2k+1)(2k+2)z^{2k}}.
 \label{eq:FL-trunc}
\end{align}
Then
\begin{equation}\label{eq:FL-expansion}
 \Log G(z+1)=\mathcal B_N(z)+R_N(z).
\end{equation}
For \(|\arg z|<\pi/2\), one exact remainder representation is
\begin{equation}\label{eq:Nemes-remainder}
 R_N(z)=\frac{(-1)^N}{2\pi^2z^{2N}}
 \int_0^\infty
 \frac{t^{2N-1}\operatorname{Li}_2(\e^{-2\pi t})}
      {1+(t/z)^2}\dd t.
\end{equation}
Equivalent periodic-Bernoulli integrals continue it to the cut plane.  They use
only a fixed Bernoulli polynomial, rational powers, and an improper integral
with exponential or algebraic tails, so they fit the framework of
\cref{sec:ipra}.

\begin{theorem}[Explicit Barnes remainder]\label{thm:Nemes-bound}
Let \(N\ge1\), \(z\ne0\), and \(\theta=\arg z\).  Then
\begin{equation}\label{eq:Nemes-bound}
 |R_N(z)|\le
 \frac{|B_{2N+2}|}
 {2N(2N+1)(2N+2)|z|^{2N}}
 \sec^{2N+1}\!\left(\frac\theta2\right),
 \qquad |\theta|<\pi.
\end{equation}
On \(|\theta|\le\pi/4\), the secant factor may be replaced by \(1\).  If
\(z>0\), \(R_N(z)\) has the sign of the first omitted term and smaller absolute
value.
\end{theorem}

\begin{proof}[Constructive reading of the remainder proof]
In \(|\theta|<\pi/2\), use \cref{eq:Nemes-remainder} and bound
\(1/|1+(t/z)^2|\) by a rational upper approximation to the appropriate angular
factor.  Expand
\(
\operatorname{Li}_2(\e^{-2\pi t})
 =\sum_{m\ge1}\e^{-2\pi mt}/m^2
\), integrate each positive term, and use the standard Bernoulli evaluation of
the resulting zeta value.  This produces exactly the coefficient in
\cref{eq:Nemes-bound}.  For wider sectors, rotate the integration path or use
the periodic-Bernoulli representation and the elementary inequality
\[
 |t+z|\ge (t+|z|)\cos(\theta/2).
\]
Raising it to the required power gives the secant factor.  On the positive
axis the integrand in a sign-normalized representation has constant sign and a
denominator strictly larger than one, proving both sign and first-omitted-term
bounds.

For a formal fixed-sector theorem, the input contains a rational
\(c>0\) with
\(
\cos(\theta/2)\ge c
\).
The verifier uses \(c^{-(2N+1)}\), avoiding an angle or inverse-trigonometric
function in the object language.  All integral cutoffs have primitive-recursive
bounds because the kernel is exponentially decaying near infinity and
polynomially regular at zero.
\end{proof}

Thus the DLMF expansion is not merely formal.  For every fixed \(N\), sector
certificate and precision, \cref{eq:Nemes-bound} gives a decidable threshold on
\(|z|\).

\subsection{The Barnes--Stirling form on the target pages}

Insert the fixed-order Stirling expansion
\begin{equation}\label{eq:stirling-gamma-use}
 \Log\Gamma(z+1)=
 \left(z+\frac12\right)\Log z-z+\frac12\log(2\pi)
 +\sum_{r=1}^{N}\frac{B_{2r}}{2r(2r-1)z^{2r-1}}
 +S_N(z)
\end{equation}
into \cref{eq:FL-expansion}.  The \(r=1\) term contributes \(1/12\).  For
\(k\ge1\), the \(r=k+1\) Gamma coefficient combines with the Barnes coefficient
as
\begin{align*}
 &\frac{B_{2k+2}}{(2k+2)(2k+1)z^{2k}}
 +\frac{B_{2k+2}}{2k(2k+1)(2k+2)z^{2k}}\\
 &\hspace{3em}=\frac{B_{2k+2}}{4k(k+1)z^{2k}}.
\end{align*}
Therefore
\begin{align}
 \Log G(z+1)={}&
 \frac{z^2}{2}\Log z-\frac{3z^2}{4}
 +\frac z2\log(2\pi)-\frac1{12}\Log z
 +\left(\frac1{12}-\log A\right)\notag\\
 &+\sum_{k=1}^{N-1}
 \frac{B_{2k+2}}{4k(k+1)z^{2k}}
 +\widetilde R_N(z),
 \label{eq:Barnes-asymptotic}
\end{align}
where
\begin{equation}\label{eq:combined-remainder}
 \widetilde R_N(z)=R_N(z)+zS_N(z).
\end{equation}
The Gamma packet supplies an explicit fixed-sector bound for \(S_N\), while
\cref{eq:Nemes-bound} bounds \(R_N\).  Hence a fully explicit bound is
\begin{align}
 |\widetilde R_N(z)|\le{}&
 \frac{|B_{2N+2}|}{|z|^{2N}}
 \left[
 \frac{C_N^{\Gamma}(\delta)}{(2N+1)(2N+2)}
 +\frac{C_N^G(\delta)}{2N(2N+1)(2N+2)}
 \right],
 \label{eq:combined-bound}
\end{align}
where \(C_N^G(\delta)\) may be taken as the rational upper bound on
\(\sec^{2N+1}(\theta/2)\), and \(C_N^\Gamma(\delta)\) is the corresponding
Stirling-remainder factor from the fixed sector
\(|\arg z|\le\pi-\delta\).  These constants are finite rational certificate
fields, not implicit big-O constants.

On the positive real axis both component remainders have the sign of their
first omitted terms and smaller magnitude.  Their first omitted coefficients
sum exactly to
\[
 \frac{|B_{2N+2}|}{4N(N+1)z^{2N}}.
\]
Thus the particularly clean bound
\begin{equation}\label{eq:positive-Barnes-bound}
 |\widetilde R_N(z)|
 <\frac{|B_{2N+2}|}{4N(N+1)z^{2N}}
 \qquad(z>0)
\end{equation}
holds, with the sign of the omitted Barnes term.  This is stronger than the
big-O statement displayed on both target pages.

Taking only the leading part of \cref{eq:Barnes-asymptotic} gives
\begin{equation}\label{eq:leading-asymptotic}
 \Log G(z+1)
 =z^2\left(\frac12\Log z-\frac34\right)
 +\frac z2\log(2\pi)-\frac1{12}\Log z+\zeta'(-1)
 +O(z^{-2}),
\end{equation}
because \(\zeta'(-1)=1/12-\log A\).  The weaker \(O(1/z)\) displayed by one
source is of course true, but the actual algebraic series contains only even
negative powers.

\subsection{An iPRA asymptotic contract}

Fix externally:
\begin{itemize}
\item a truncation order \(N\ge1\);
\item a rational sector-gap certificate \(c\le\cos(\arg z/2)\), \(c>0\);
\item a requested precision \(p\);
\item a rational lower bound \(|z|\ge L\).
\end{itemize}
The verifier computes the rational right side of
\cref{eq:Nemes-bound} or \cref{eq:combined-bound} and checks that it is at most
\(2^{-p}\) after clearing denominators.  It then checks finite balls for every
term of the truncation.  This is a quantifier-free predicate.  A producer can
increase \(L\) by bounded search until the inequality holds.  The theorem
``the asymptotic expansion is valid'' is the schema asserting soundness of this
procedure for every \(p\), not a formula containing a primitive big-O symbol.

\subsection{Order two and infinite type}

Let
\(
M_G(r)=\max_{|z|\le r}|G(z)|
\)
in classical notation.  Constructively, we need only an upper bound valid for
all named inputs in the disk and a single lower-witness input.

\begin{lemma}[Elementary global upper bound]\label{lem:order-upper}
There is an explicit integer constant \(C\) such that for every real
\(r\ge2\) and \(|z|\le r\),
\begin{equation}\label{eq:order-upper}
 \log|G(1+z)|\le C r^2\log(r+2).
\end{equation}
One may take a deliberately coarse value such as \(C=12\).
\end{lemma}

\begin{proof}
Split the canonical product at \(2r\).  For \(k\le2r\),
\[
 \log|E_k(z)|
 \le k\log(1+r/k)+r+\frac{r^2}{2k}.
\]
Sum the first term using
\(\log(1+r/k)\le\log(r+1)\) and
\(\sum_{k\le2r}k\le3r^2\).  The second contributes at most \(2r^2\), and the
third at most
\(
\frac12r^2(1+\log(2r))
\).
The logarithmic tail beyond \(2r\) is at most \(r^2/3\) by
\cref{eq:qtail}.  The prefactor is bounded by a fixed multiple of \(r^2\).
Collecting the terms gives \cref{eq:order-upper}.  Every real inequality used
has a rational majorant version.
\end{proof}

\begin{lemma}[Elementary integer lower bound]\label{lem:order-lower}
For all sufficiently large integers \(n\),
\begin{equation}\label{eq:order-lower}
 \log G(n+1)\ge \frac{n^2}{20}\log(n/2).
\end{equation}
\end{lemma}

\begin{proof}
From the exact integer formula,
\[
 \log G(n+1)=\sum_{j=1}^{n-1}(n-j)\log j.
\]
Restrict to integers \(j\) with \(n/2\le j\le3n/4\).  There are at least
\(n/5\) of them once \(n\) exceeds a fixed numeral, each satisfies
\(n-j\ge n/4\), and \(\log j\ge\log(n/2)\).  Their contribution proves
\cref{eq:order-lower}.  Floors and the small exceptional values are handled by
bounded numeral checks.
\end{proof}

\begin{theorem}[Order and type]\label{thm:order-type}
The Barnes $G$-function is entire of order exactly two and infinite type at
order two.
\end{theorem}

\begin{proof}
The upper bound implies
\[
 \limsup_{r\to\infty}\frac{\log\log M_G(r)}{\log r}\le2.
\]
The lower bound, using the witness \(z=n+1\), gives the reverse inequality.
Moreover,
\[
 \frac{\log M_G(n+1)}{(n+1)^2}
 \ge c\log(n/2)\longrightarrow\infty,
\]
so the type at order two is infinite.

The constructive replacement is explicit.  Given a positive rational
\(\varepsilon\), bounded search finds a threshold after which
\(C r^2\log(r+2)\le r^{2+\varepsilon}\); the lower witness exceeds
\(r^{2-\varepsilon}\) after another threshold.  Given an integer \(T\), choose
an explicit power of two \(n=2^{a(T+1)}\), with fixed \(a\) large enough that a
rational lower bound for \(\log2\) proves
\(\frac1{20}\log(n/2)>T\).  Then
\(\log|G(n+1)|>T(n+1)^2\) after increasing \(a\) to absorb the harmless ratio.
All thresholds are primitive recursive.
\end{proof}

\subsection{What is and is not formalized here}

The algebraic asymptotic expansion and the explicit remainder contracts are
fully developed mathematically.  The package does not claim a completed Lean
port of Nemes' periodic-Bernoulli integral.  The formalization plan is direct:
fixed Bernoulli polynomials, a rational sector lower bound, finite quadrature,
and algebraic tail integration.  Exponentially improved Stokes expansions are
outside the two target pages and are not needed for any theorem in the coverage
ledger.
